nano tokenize_data.py   # paste, save

python -c "import ast;ast.parse(open('tokenize_data.py').read());print('OK')"

python tokenize_data.py

git add notes.tex tokenize_data.py && git commit -m "Add notes and tokenizer application" && git push\documentclass[11pt]{article}

\usepackage[utf8]{inputenc}

\usepackage[margin=1in]{geometry}

\usepackage{booktabs}

\usepackage{parskip}

\usepackage{hyperref}

\usepackage{enumitem}


\title{Project 1 --- Running Notes\\[0.4em]

\large Small Code Language Model, Chinchilla-Optimal Rebuild\\

CS 5090 --- Software Engineering with Generative AI}

\author{Marjan Beheshti}

\date{Last updated: October 22, 2026}


\begin{document}

\maketitle


\emph{Working log. Every decision is recorded with the reason and the evidence

behind it, so the final report can be assembled from this file.}


\tableofcontents

\newpage


% ============================================================

\section{Why this rebuild exists}

% ============================================================


A first attempt trained a 16.9M-parameter LLaMA-style model on a

security-focused Python corpus ($\approx$0.85--4.9M tokens). A controlled

data-scaling study (D1/D2/D3) showed best validation loss falling

$3.12 \to 2.71 \to 2.28$ while the \texttt{ast.parse} rate of generated code

stayed at or below 4\% at every data size.


The diagnosis was that the model sat roughly \textbf{400$\times$ below}

Chinchilla-optimal data ($\approx$0.06--0.36 tokens/parameter against a target

of $\approx$20). That left an open question the old setup could not answer:

\emph{does code validity emerge once the model is actually trained on enough

data, or is 16.9M parameters simply too small?}


\textbf{Key realisation:} the binding constraint was never compute. At

$\approx$32{,}768 tokens per step, 340M tokens is roughly 10{,}000 steps, which

is 1--2 GPU-hours on this hardware. The constraint was corpus size. This rebuild

therefore targets $\approx$20 tokens/parameter on a general Python corpus.


% ============================================================

\section{Infrastructure}

% ============================================================


\begin{center}

\begin{tabular}{ll}

\toprule

\textbf{Component} & \textbf{Detail} \\

\midrule

Platform    & Jetstream2 (NSF cloud, Indiana University) \\

Allocation  & CIS262006, CS 5090 classroom allocation \\

Instance    & \texttt{ktgap-gpu}, Ubuntu 22 \\

GPU         & NVIDIA A100 slice (\texttt{GRID A100X-20C}, 20\,GB) \\

CPU / RAM   & 16 cores / 60\,GB \\

Storage     & volume \texttt{ktgap-data}, 98\,GB (44\,GB free at start) \\

PyTorch     & 2.11.0+cu128, CUDA available \\

Repository  & \url{https://github.com/marjanbeheshti/cs5090-Project-1} \\

Working dir & \texttt{/media/volume/ktgap-data/cs5090-p1} \\

\bottomrule

\end{tabular}

\end{center}


A project-local \texttt{venv} is used rather than the existing \texttt{ktgap}

environment, so this project cannot disturb unrelated work on the same server.


\textbf{Cost discipline:} the instance bills 32 SU/hour on \texttt{g3.large}

while running, even when idle, against a 20{,}000 SU pool shared by the whole

class. Shelving between sessions is mandatory; the volume retains all data.


% ============================================================

\section{Dataset selection}

% ============================================================


Two candidates were evaluated by inspecting raw records before any pipeline was

written. This inspection-first step was decisive.


\subsection{Rejected: \texttt{h4iku/coconut\_python2010}}

Chosen initially because it ships buggy/fixed pairs with a \texttt{context}

column described as the inlined buggy function. Inspection showed the context

is \textbf{whitespace-flattened}: no newlines, no indentation, with

\texttt{def f(): """doc""" assert x if flatten: ...} collapsed onto one line.

Python indentation is semantic, so such text can never pass \texttt{ast.parse}

and could not support a parse-rate metric. A second defect: sampled records had

\texttt{rem} identical to \texttt{add}, i.e.\ docstring reflows labelled as bug

fixes. \textbf{Rejected.}


\subsection{Selected: \texttt{codeparrot/codeparrot-clean}}

Real Python with genuine newlines and four-space indentation. Each record also

carries \texttt{license}, \texttt{repo\_name}, \texttt{path} and \texttt{hash},

which directly supports the data-sheet and licensing requirements of the

assignment.


% ============================================================

\section{Step 3 --- Corpus construction (\texttt{build\_corpus.py})}

% ============================================================


Filters applied in order: permissive licence only; drop autogenerated; keep

200--100{,}000 characters; \textbf{drop anything failing \texttt{ast.parse}};

drop normalised-hash duplicates; stop at the token budget.


Normalisation for deduplication collapses runs of spaces/tabs and removes blank

lines, so two files differing only in formatting hash identically. The

\emph{original} text is what is written out; the normalised form is never

training data.


\begin{center}

\begin{tabular}{lrr}

\toprule

\textbf{Outcome} & \textbf{Files} & \textbf{Share of scanned} \\

\midrule

Scanned                  & 272{,}263 & --- \\

Kept                     & 139{,}424 & 51.2\% \\

Dropped: non-permissive licence & 110{,}063 & 40.4\% \\

Dropped: unparseable     &  20{,}935 & 7.7\% \\

Dropped: too large       &   1{,}186 & 0.4\% \\

Dropped: duplicate       &     655 & 0.2\% \\

\bottomrule

\end{tabular}

\end{center}


Kept 1{,}190{,}013{,}075 characters. Licence mix of the kept set: apache-2.0

(54{,}943), mit (42{,}615), bsd-3-clause (34{,}810), bsd-2-clause (4{,}296),

unlicense (1{,}925), isc (460), cc0-1.0 (375). GPL/AGPL/LGPL were excluded.


\textbf{Note for the report:} the 7.7\% unparseable rate is largely Python 2

source (\texttt{print "x"}), since \texttt{codeparrot-clean} predates the Python 3

transition in much of its crawl. Dropping it is intentional --- the parse-rate

metric requires a Python 3 corpus.


Runtime: approximately 6 minutes, network-bound.


% ============================================================

\section{Step 4 --- Train/validation split (\texttt{split\_corpus.py})}

% ============================================================


\textbf{Repository-level holdout, not random files.} Files within one repository

share helpers, imports and style; a random file-level split lets the model train

on part of a project and be validated on the rest, which flatters validation

loss. Holding out whole repositories means a validated project was never seen at

all.


A second pass then drops any training file whose normalised hash also appears in

validation, catching files copied \emph{between} repositories, which a

repository split alone cannot catch.


\begin{center}

\begin{tabular}{lr}

\toprule

\textbf{Quantity} & \textbf{Value} \\

\midrule

Seed                          & 1337 \\

Total repositories            & 50{,}789 \\

Train repositories / files    & 50{,}455 / 138{,}631 \\

Validation repositories / files & 334 / 793 \\

Train characters              & 1{,}182{,}924{,}896 \\

Validation characters         & 7{,}088{,}179 \\

Train tokens (estimated)      & $\approx$338M \\

Cross-repository duplicates dropped & 0 \\

\bottomrule

\end{tabular}

\end{center}


Zero cross-repository duplicates were found --- \texttt{codeparrot-clean} is

already deduplicated upstream. The check is retained because ``we verified and

found none'' is a stronger claim than not checking.


The corpus is wide rather than deep: 50{,}789 repositories for 139{,}424 files,

under three files per repository on average, which should help generalisation.


Files are separated by a \texttt{<|file|>} document boundary so the model does

not read one file as running directly into the next.


% ============================================================

\section{Step 5 --- Tokenizer (\texttt{train\_tokenizer.py})}

% ============================================================


SentencePiece BPE, vocabulary 8{,}192, trained \textbf{only} on

\texttt{data/train.txt}. Letting the tokenizer see validation text would leak

information into the vocabulary and make validation loss unreliable.


\begin{center}

\begin{tabular}{lp{9cm}}

\toprule

\textbf{Setting} & \textbf{Why} \\

\midrule

\texttt{model\_type="bpe"} & Subword. Identifiers are unbounded, so word-level

would hit \texttt{<unk>} constantly; character-level would leave a 512-token

context holding barely one function. \\

\texttt{character\_coverage=1.0} & Natural-language defaults (0.9995) discard

rare characters; code cannot afford that. \\

\texttt{remove\_extra\_whitespaces=False} & \textbf{Critical.} The default

collapses whitespace, which would destroy Python indentation. \\

\texttt{normalization\_rule\_name="identity"} & Same reason --- no Unicode

normalisation applied to source text. \\

\texttt{byte\_fallback=True} & Unknown characters become raw bytes rather than

\texttt{<unk>}, so nothing is silently lost. \\

\texttt{input\_sentence\_size=2{,}000{,}000} & 1.2\,GB will not fit in memory;

lines are sampled with shuffling under a fixed seed. \\

\bottomrule

\end{tabular}

\end{center}


Special ids follow the LLaMA convention: \texttt{unk}=0, \texttt{bos}=1,

\texttt{eos}=2, \texttt{pad}=$-1$, with \texttt{<|file|>}=3.


\textbf{Verification:} a round-trip test encodes an indented function and

requires the decode to match the original exactly. It passed, confirming

indentation and newlines survive tokenization. The script exits with an error if

it ever fails, rather than allowing training on a broken tokenizer.


The \texttt{.vocab} escaped-character warnings affect only the human-readable

vocabulary dump, not the \texttt{.model} file used for encoding.


\textbf{Open number:} the round-trip sample gave 1.81 characters/token, against

the 3.5 assumed when sizing the corpus. The sample was short and

indentation-heavy, so the true ratio over full files will be higher. Step 6

measures it exactly; if the real count exceeds 338M tokens, the run is simply

above Chinchilla-optimal rather than below, which is not a problem.


% ============================================================

\section{Planned configuration}

% ============================================================


\begin{center}

\begin{tabular}{ll}

\toprule

\textbf{Setting} & \textbf{Value} \\

\midrule

Layers / heads / embedding & 6 / 6 / 384 \\

Context length             & 512 \\

Vocabulary                 & 8{,}192 \\

Parameters                 & $\approx$16.9M \\

Target tokens              & $\approx$338M ($\approx$20 tokens/parameter) \\

Precision                  & bf16 \\

\bottomrule

\end{tabular}

\end{center}


Configuration is held identical to the earlier security-corpus runs, so the two

projects remain comparable and the only changed variable is the data.


% ============================================================

\section{Open items}

% ============================================================


\begin{itemize}[nosep]

  \item \textbf{Step 6} --- tokenize to \texttt{train.bin}/\texttt{val.bin};

        obtain the true token count.

  \item \textbf{Step 7} --- baseline training run.

  \item \textbf{Ablation} (Phase 1, graded) --- not yet run. Prediction must be

        recorded \emph{before} the run. Character-level tokenizer is the natural

        choice given the BPE rationale above.

  \item \textbf{Phase 2} (graded) --- 20 real inputs with accept / accept-after-edit

        / reject counts, plus best, worst and near-miss examples.

  \item \textbf{Phase 3} (graded) --- 20--30 probes, 10+ documented failures in a

        CSV using the assignment's exact category names.

  \item \textbf{Phase 4} --- intervention and re-run of the same probes.

\end{itemize}


\subsection{Bug-injection plan (decide after the baseline)}

Rather than a mined bug-fix dataset, single-statement bugs are injected into

held-out clean functions by AST transformation (flip a comparison operator, swap

two variable names, change a literal by one). This yields real parseable Python

on both sides with exact ground truth, which neither candidate dataset provided.


\begin{itemize}[nosep]

  \item \textbf{Option A (evaluation only).} Use injected bugs as the Phase 3

        probe set. No extra training. Low cost.

  \item \textbf{Option B (adds fine-tuning).} Also fine-tune on

        buggy$\rightarrow$fixed pairs as the Phase 4 intervention, scored by

        exact match on the fixed line. Exact-match scoring sidesteps the

        parse-rate floor that made the previous project's D1/D2/D3 runs

        indistinguishable.

\end{itemize}


Given the deadline, Option A is the default; Option B is added only if the

baseline completes with time to spare.


\end{document}
